\begin{lemma}[Scheide enlarged set is closed (19)]
The enlarged set $X=V(T)\cup\bigcup_\delta U_\delta$ is closed with
respect to $\varphi$.
\end{lemma}
\begin{proof}
If there is no defective color, the indexed union in
\noderef{ScheideEnlargedVertexSet} is empty, so $X=V(T)$, which is closed
by the globally maximum clause and \noderef{ScheideGlobalMaximumOrder}.
Assume henceforth that a defective color exists.

Suppose an edge identity $f=uw$ of color $\alpha$ leaves $X$, while
$\alpha$ is missing at some vertex of $X$.  If $u$ lies in a replacement
pair, choose its index $\delta$; if $u\in V(T)$, choose any defective
color, which exists by the preceding case split.  In either event $u$
belongs to $X_\delta=V(T)\cup U_\delta$.  By
\noderef{ScheideReplacementUnionClosed}, $X_\delta$ is closed, so the
vertex missing $\alpha$ is not in $X_\delta$.  By the equal-or-disjoint
alternative \noderef{ScheideReplacementPairsLaminar}, it belongs to a
pair $U_{\delta'}$ disjoint from $U_\delta$.  Closedness of
$X_{\delta'}$ then shows that $u$ cannot lie in $V(T)$; hence
$u\in U_\delta$.

Consider the $(\alpha_0,\alpha)$ incidence component through $u$.  Both
colors are missing on the closed set $X_{\delta'}$, so the component avoids
that set, in particular $V(T)$.  By elementariness of $X$ from
\noderef{ScheideEnlargedElementary}, $\alpha$ is not a block color, since
every block color already has a missing vertex in $V(T)$.  Also
$\alpha\ne\delta$: the named $\delta$-edge from $V(T)$ to $U_\delta$
crosses the boundary of the closed set $X_{\delta'}$, so $\delta$ is
present throughout that set.  The setup gives
$\alpha_0\notin\Gamma$, and $\alpha_0\ne\delta$ because $\delta$ is
present on the original closed tree.

The component is therefore identical under $\varphi$ and
$\varphi_\delta$, and it contains no edge of the replacement tree
$T_\delta$: its inherited edges have an end or missing witness in the
avoided original tree, while its two new edges have color $\delta$.
Switching this full incidence component preserves properness by
\noderef{TwoColorComponentSwapProper} and fixes the unique uncolored edge.
Every replacement-tree edge and every missing witness for it remains
unchanged, so \noderef{ScheideReplacementMaximum} remains a globally
maximum Tashkinov triple.  The boundary edge $f$ lies in the switched
component and acquires color $\alpha_0$, while $\alpha_0$ remains missing
at $y_0$.  Appending $f,w$ gives a larger Tashkinov tree, contradicting
global maximality.  Thus no such boundary edge exists, exactly the closure
condition of \noderef{ColoredVertexSetClosed}.
\end{proof}
